\section{Discussion and Research Agenda}\label{sec:discussion}

\textbf{Measurement changes the maintenance story.} Under Qwen, 10\% of initially supported claims receive a detection by touch four, and roughly 15\% by touch six. A full Llama replay of the same 129 functions gives $\hat S(6)=\XLlamaFullSsix$ [\XLlamaFullSsixLo,\XLlamaFullSsixHi] among \XLlamaFullN{} Llama-supported anchor claims, compared with \XMainN{} under Qwen. Anchor agreement is \XLlamaAnchorAgree\% ($\kappa=\XLlamaAnchorKappa$; 1,725 labels). The replay changes both eligibility and subsequent judgments, so it cannot isolate follow-up error alone. The gap would nevertheless materially affect a refresh schedule; choosing one from the primary curve would outrun the validation evidence.

\textbf{Human review changes the event count.} An independent rater relabeled all \XHumanEventAuditN{} unique observed Qwen-CON rows through touch six: \XHumanLabelConN{} contradicted, \XHumanLabelSupN{} supported and \XHumanLabelUnrN{} unresolved. Confirmation among resolved rows is \XHumanPpv\% [\XHumanPpvLo,\XHumanPpvHi\% Wilson]. Substitution yields $\hat S(6)=\XHumanLabelSsixLo$ when unresolved is an event or \XHumanLabelSsixHi{} when censored (function-cluster bootstrap \XHumanSsixBootLo--\XHumanSsixBootHi). It removes observed false positives but cannot recover missed contradictions, including after a rejected first event in the alive-only archive. The increase demonstrates the consequence of event verification; it does not establish a corrected semantic curve.

The two-rater audit and event review answer different questions. A1/A2 agree on \XTwoRaterAgree\% of \XTwoRaterN{} records ($\kappa=\XTwoRaterKappa$), yet the event rater confirms only \XHumanOverlapCon{} of \XHumanOverlapN{} shared Qwen-CON rows. Agreement within one frame therefore does not establish calibration across raters and frames. Independence and blinding are self-reported; no calibration session is documented. A next audit should use common guidelines and independently adjudicate both flagged and unflagged claims.

\textbf{Controlled checks have limited reach.} A synthetic suite detects \XMutationRecallPct\% of 100 known changes [\XMutationRecallLo--\XMutationRecallHi\% Wilson] with \XMutationFprPct\% false positives among 100 unchanged controls. Yet targeted review supports 75/77 Qwen-SUP/Llama-CON anchor rows and contradicts only 9/37 Qwen-CON/Llama-SUP rows. Success on controlled mutations can coexist with ambiguity in repository histories. These selected checks do not estimate population precision or recall. Likewise, churn offers a candidate review priority, but moderate discrimination and dependence on Qwen labels leave operational value open.

\textbf{What to validate next.} First, compare published JIT detectors and execution-based checks on frozen eligible inputs, reporting coverage as well as errors. Second, independently validate tracking across refactorings and renames: false matches create spurious follow-up, while failed matches may censor informatively. Third, collect representative adjudication of flagged and unflagged claims before fitting contamination-aware survival methods~\cite{bladt2026contaminated}; the present data do not identify that error mechanism. Future-conditioned selection, sparse random-cohort events, \XNRepos{} Python projects, one generator/prompt and non-local dependency changes constrain generalization. A larger, prospectively fixed cohort is needed before testing refresh policies or developer outcomes.
